%\title{Hierarchical Auxetic Elasticity}

%\author{Carlos Villarroel${}^{1}$ and  Gustavo D"uring${}^{1}$}
%\affiliation{${}^1$Instituto de F\'isica, Pontificia Universidad Cat\'olica de Chile, Casilla 306, Santiago, Chile}

\documentclass[11pt,a4paper]{article}

% -----------------------------
% Idioma y codificación
% -----------------------------
\usepackage[spanish]{babel}
\usepackage[utf8]{inputenc}
\usepackage[T1]{fontenc}

% -----------------------------
% Matemática
% -----------------------------
\usepackage{amsmath,amssymb}

% -----------------------------
% Figuras
% -----------------------------
\usepackage{graphicx}
\usepackage{float}
\usepackage[font=footnotesize,labelfont=bf]{caption}

% -----------------------------
% Márgenes
% -----------------------------
\usepackage{geometry}
\geometry{
  left=2.5cm,
  right=2.5cm,
  top=2.5cm,
  bottom=2.5cm
}

% -----------------------------
% Espaciado
% -----------------------------
\usepackage{setspace}
\setstretch{1.15}

% -----------------------------
% Bibliografía
% -----------------------------
\usepackage[numbers]{natbib}

% -----------------------------
% Información del documento
% -----------------------------
\title{\textbf{Hierarchical Auxetic Elasticity}}
\author{
Carlos Villarroel$^{1}$ and Gustavo D\"uring$^{1}$ \\
\vspace{0.2cm}
\small $^{1}$Instituto de Física, Pontificia Universidad Católica de Chile, Santiago, Chile
}
\date{Junio 2025}

% -----------------------------
\begin{document}
% -----------------------------

\maketitle

% =================================================
\section{Modelo}
\label{sec:model}
% =================================================

\begin{figure}[H]
  \centering
  \includegraphics[width=0.45\textwidth]{figures/fractal1.png}
  \caption{Estructura fractal con niveles 0, 1 y 2, mostrando el crecimiento jerárquico. Cada nivel es una replicación del anterior. El nivel 1 está compuesto por cuatro cuadrados que pueden rotar alrededor de vértices internos y externos.}
  \label{fig:levels}
\end{figure}

Siguiendo la Fig.~\ref{fig:levels}, definimos la energía potencial de una estructura
jerárquica con $N$ niveles como
\begin{equation}
E_N = \sum_{i=1}^N 4^{N-i}\,E^{(i)},
\label{eq:En_sum}
\end{equation}
donde el término correspondiente al nivel $i$ está dado por
\begin{equation}
E^{(i)} = 2k_i\Big[(\phi_i^+ - \phi_{i,0}^+)^2 + (\phi_i^- - \phi_{i,0}^-)^2\Big].
\end{equation}

Aquí $k_i$ es la constante elástica asociada al nivel $i$, y $\phi_i^\pm$ son los
ángulos internos de rotación definidos en la Fig.~\ref{fig:angles}. 
A diferencia del caso idealizado donde el estado no deformado corresponde a
$\phi_{i,0}^+=\pi$ y $\phi_{i,0}^-=0$, consideramos un estado inicial más general,
caracterizado por un conjunto de ángulos $\{\delta_i\}$ que describen desviaciones
respecto a esa configuración.

Introducimos un único ángulo elemental por nivel, $\theta_i$, mediante las relaciones
\begin{equation}
\phi_i^+ = \pi - (\theta_i + S_i),
\qquad
\phi_i^- = \theta_i - S_i,
\label{eq:phi_combined}
\end{equation}
donde
\begin{equation}
S_i = \sum_{j=1}^{i-1}\theta_j
\end{equation}
representa la rotación acumulada inducida por los niveles superiores.

El estado inicial se parametriza entonces mediante
\begin{equation}
\phi_{i,0}^+ = \pi - (\delta_i + R_i),
\qquad
\phi_{i,0}^- = \delta_i - R_i,
\end{equation}
donde
\begin{equation}
R_i = \sum_{j=1}^{i-1}\delta_j
\end{equation}
es la rotación acumulada asociada al estado inicial.

\begin{figure}[H]
  \centering
  \includegraphics[width=0.45\textwidth]{figures/fractal2.png}
  \caption{Definición de los ángulos $\phi_i^\pm$ y su descomposición en términos de los ángulos elementales $\theta_i$.}
  \label{fig:angles}
\end{figure}

Sustituyendo estas relaciones en la Ec.~\eqref{eq:En_sum} y midiendo la energía
respecto al estado inicial, se obtiene una expresión completamente en función
de los ángulos $\{\theta_i\}$:
\begin{equation}
E_N = \sum_{i=1}^{N} 4^{N-i+1} k_i
\left[
(\theta_i - \delta_i)^2 + (S_i - R_i)^2
\right].
\label{eq:En_theta_final}
\end{equation}

Esta forma muestra explícitamente que la energía elástica depende únicamente
de las desviaciones de los ángulos $\theta_i$ respecto al estado inicial
$\delta_i$, así como de las diferencias entre las rotaciones acumuladas
$S_i$ y $R_i$.


% =================================================
\section{Tamaños del sistema}
\label{sec:size}
% =================================================

Para caracterizar la extensión geométrica de la estructura se definen cuatro tamaños
característicos: la extensión máxima y mínima en las direcciones horizontal y
vertical, denotadas respectivamente por
$Lx_{\max}$, $Lx_{\min}$, $Ly_{\max}$ y $Ly_{\min}$.

En el nivel base ($i=0$), estas cantidades satisfacen las condiciones iniciales
\begin{equation}
Lx_{\max}^{(0)} = Lx_{\min}^{(0)} =
Ly_{\max}^{(0)} = Ly_{\min}^{(0)} = 2^{-N},
\end{equation}
de modo que el tamaño global del sistema en el estado inicial sea del orden de
la unidad.

\begin{figure}[H]
  \centering
  \includegraphics[width=0.45\textwidth]{figures/fractal3.png}
  \caption{Definición de los tamaños característicos del sistema en las direcciones horizontal y vertical.}
  \label{fig:size}
\end{figure}

La construcción jerárquica implica que estos tamaños se actualizan de manera iterativa
al pasar de un nivel al siguiente. Para cada nivel $i$, las relaciones de recurrencia
son
\begin{align}
Ly_{\max}^{(i)} &=
2\sin\!\left(\tfrac{\theta_i}{2}\right)Lx_{\max}^{(i-1)}
+
2\cos\!\left(\tfrac{\theta_i}{2}\right)Ly_{\min}^{(i-1)}, \\
Lx_{\min}^{(i)} &=
2\sin\!\left(\tfrac{\theta_i}{2}\right)Ly_{\min}^{(i-1)}, \\
Lx_{\max}^{(i)} &=
2\sin\!\left(\tfrac{\theta_i}{2}\right)Ly_{\min}^{(i-1)}
+
2\cos\!\left(\tfrac{\theta_i}{2}\right)Lx_{\max}^{(i-1)}, \\
Ly_{\min}^{(i)} &=
2\cos\!\left(\tfrac{\theta_i}{2}\right)Ly_{\min}^{(i-1)}.
\end{align}

Estas relaciones dependen únicamente de los ángulos geométricos $\theta_i$ que
describen la configuración instantánea del sistema. En particular, aunque en la
sección anterior introdujimos un estado inicial caracterizado por los ángulos
$\{\delta_i\}$ para definir la energía elástica, la geometría del sistema está
completamente determinada por los valores actuales de $\{\theta_i\}$.

En particular, la extensión horizontal máxima del sistema al nivel $N$ puede
escribirse de forma cerrada como
\begin{equation}
Lx_{\max}^{(N)} =
\prod_{i=1}^{N}
\left[
\sin\!\left(\tfrac{\theta_i}{2}\right)
+
\cos\!\left(\tfrac{\theta_i}{2}\right)
\right].
\label{eq:Lxmax_product}
\end{equation}

% =================================================
\section{Hamiltoniano y condiciones de equilibrio}
\label{sec:hamiltonian}
% =================================================

Para estudiar la relajación mecánica de la estructura jerárquica utilizamos como
parámetro de control la extensión horizontal máxima del sistema. A partir de esta
sección, definimos
\begin{equation}
L \equiv Lx_{\max}^{(N)}(\{\theta_i\}),
\end{equation}
entendida como una función de los ángulos elementales $\{\theta_i\}$ que caracteriza
la geometría global del sistema.

La configuración mecánica bajo una extensión horizontal impuesta $\ell$ se describe
mediante un Hamiltoniano efectivo que combina la energía elástica interna con un
término de Lagrange. En términos de los ángulos $\{\theta_i\}$, el Hamiltoniano se
escribe como
\begin{equation}
\mathcal{H}(\{\theta_i\}) =
\sum_{i=1}^{N} 4^{N-i+1} k_i
\left[
(\theta_i-\delta_i)^2 + (S_i-R_i)^2
\right]
-
f\left[L(\{\theta_i\}) - \ell\right],
\label{eq:H_def_ki}
\end{equation}
donde $f$ es la fuerza efectiva conjugada a la extensión horizontal, mientras que
$\{\delta_i\}$ caracterizan el estado inicial de referencia introducido en la
sección anterior. En esta expresión
\begin{equation}
S_i=\sum_{j=1}^{i-1}\theta_j,
\qquad
R_i=\sum_{j=1}^{i-1}\delta_j .
\end{equation}

En el modelo geométrico considerado, la extensión $L$ está dada explícitamente por
\begin{equation}
L(\{\theta_i\}) =
\prod_{i=1}^{N}
\left(
\sin\!\tfrac{\theta_i}{2}
+
\cos\!\tfrac{\theta_i}{2}
\right).
\label{eq:L_def}
\end{equation}

El estado de equilibrio del sistema se obtiene imponiendo que el Hamiltoniano sea
estacionario frente a variaciones independientes de todos los ángulos, es decir,
\begin{equation}
\frac{\partial \mathcal{H}}{\partial \theta_j}=0,
\qquad
j=1,\dots,N.
\label{eq:stationary}
\end{equation}

\subsection*{Condición de equilibrio para el nivel superior}

El ángulo correspondiente al último nivel, $\theta_N$, juega un rol particular ya que
no contribuye a ningún término acumulado $S_i$ con $i>N$. En consecuencia, la derivada
del Hamiltoniano con respecto a $\theta_N$ se simplifica. Derivando la
Ec.~\eqref{eq:H_def_ki} se obtiene
\begin{equation}
\boxed{
\frac{\partial \mathcal{H}}{\partial \theta_N}
=
2\cdot 4^{1} k_N\,(\theta_N-\delta_N)
-
\frac{f}{2}\,L\,
\frac{\cos\!\left(\tfrac{\theta_N}{2}\right)-\sin\!\left(\tfrac{\theta_N}{2}\right)}
{\cos\!\left(\tfrac{\theta_N}{2}\right)+\sin\!\left(\tfrac{\theta_N}{2}\right)}
=0.}
\label{eq:dH_thetaN_ki}
\end{equation}



% =================================================
\section{Aproximación de pequeños ángulos}
\label{sec:small}
% =================================================

En esta sección analizamos el régimen de pequeños ángulos, donde tanto los ángulos
geométricos $\theta_i$ como los parámetros del estado inicial $\delta_i$ satisfacen
$\theta_i,\delta_i\ll 1$, y mantenemos términos hasta segundo orden. Recordemos que
el Hamiltoniano está dado por la Ec.~\eqref{eq:H_def_ki}, con la extensión geométrica
$L(\{\theta_i\})$ definida en la Ec.~\eqref{eq:L_def}.

\subsection{Expansión de la extensión geométrica}

Para $\theta_i\ll 1$ se cumple
\begin{equation}
\sin\!\left(\tfrac{\theta_i}{2}\right)+\cos\!\left(\tfrac{\theta_i}{2}\right)
\simeq
1+\tfrac{\theta_i}{2}-\tfrac{\theta_i^2}{8},
\end{equation}
y por lo tanto
\begin{equation}
L(\{\theta_i\})
=
\prod_{i=1}^{N}\left(1+\tfrac{\theta_i}{2}-\tfrac{\theta_i^2}{8}\right)
\simeq
\exp\!\left(\tfrac12\sum_{i=1}^{N}\theta_i-\tfrac14\sum_{i=1}^{N}\theta_i^2\right).
\label{eq:L_small}
\end{equation}

Definimos, para abreviar, el largo aproximado
\begin{equation}
L_{\text{quad}}(\{\theta_i\})
:=
\exp\!\left(\tfrac12\sum_{i=1}^{N}\theta_i-\tfrac14\sum_{i=1}^{N}\theta_i^2\right).
\end{equation}

\subsection{Hamiltoniano truncado a segundo orden}

Bajo esta aproximación, el Hamiltoniano toma la forma
\begin{equation}
\mathcal{H}_{\text{quad}}(\{\theta_i\})
=
\sum_{i=1}^{N}4^{N-i+1}k_i\left[(\theta_i-\delta_i)^2+(S_i-R_i)^2\right]
-
f\left[L_{\text{quad}}(\{\theta_i\})-\ell\right],
\label{eq:H_quad}
\end{equation}
donde
\begin{equation}
S_i=\sum_{m=1}^{i-1}\theta_m,
\qquad
R_i=\sum_{m=1}^{i-1}\delta_m,
\qquad (S_1=R_1=0).
\end{equation}

La condición de equilibrio se obtiene imponiendo estacionariedad,
\begin{equation}
\frac{\partial\mathcal{H}_{\text{quad}}}{\partial\theta_j}=0,
\qquad j=1,\dots,N.
\label{eq:eq_cond_quad}
\end{equation}

\subsection{Derivadas parciales}

Ya que $\partial S_i/\partial\theta_j=1$ para $i>j$ y $0$ en caso contrario, la derivada
de la parte elástica puede escribirse como
\begin{equation}
\frac{\partial}{\partial\theta_j}
\sum_{i=1}^{N}4^{N-i+1}k_i\left[(\theta_i-\delta_i)^2+(S_i-R_i)^2\right]
=
2\cdot4^{N-j+1}k_j\,(\theta_j-\delta_j)
+
2\sum_{i=j+1}^{N}4^{N-i+1}k_i\,(S_i-R_i).
\label{eq:dE_dtheta_general_ki}
\end{equation}

Por otro lado, usando $L_{\text{quad}}=\exp(A)$ con
$A=\frac12\sum_i\theta_i-\frac14\sum_i\theta_i^2$, se obtiene
\begin{equation}
\frac{\partial L_{\text{quad}}}{\partial\theta_j}
=
L_{\text{quad}}\left(\tfrac12-\tfrac12\theta_j\right)
=
\frac{L_{\text{quad}}}{2}(1-\theta_j).
\label{eq:dLquad_dtheta}
\end{equation}

Combinando \eqref{eq:dE_dtheta_general_ki} y \eqref{eq:dLquad_dtheta}, la condición
\eqref{eq:eq_cond_quad} conduce a
\begin{equation}
\frac{\partial\mathcal{H}_{\text{quad}}}{\partial\theta_j}
=
2\cdot4^{N-j+1}k_j\,(\theta_j-\delta_j)
+
2\sum_{i=j+1}^{N}4^{N-i+1}k_i\,(S_i-R_i)
-
\frac{f}{2}\,L_{\text{quad}}(1-\theta_j)
=0.
\label{eq:dH_dtheta_ki}
\end{equation}

\subsection{Estructura recursiva}

Para analizar la estructura de soluciones, es útil considerar la diferencia
\begin{equation}
\Delta_j
:=
\frac{\partial\mathcal{H}_{\text{quad}}}{\partial\theta_{j+1}}
-
\frac{\partial\mathcal{H}_{\text{quad}}}{\partial\theta_{j}},
\qquad j=1,\dots,N-1.
\end{equation}

Al restar las ecuaciones \eqref{eq:dH_dtheta_ki} consecutivas, el término constante
$-\frac{f}{2}L_{\text{quad}}$ se cancela. Además, la diferencia de las sumas ponderadas
colapsa a un término local proporcional a $k_{j+1}(S_{j+1}-R_{j+1})$, obteniéndose
\begin{equation}
\Delta_j
=
2\cdot4^{N-j}\left[k_{j+1}(\theta_{j+1}-\delta_{j+1})-4k_j(\theta_j-\delta_j)-k_{j+1}(S_{j+1}-R_{j+1})\right]
+
\frac{f}{2}L_{\text{quad}}(\theta_{j+1}-\theta_j)
=0,
\label{eq:delta_ki}
\end{equation}
donde $S_{j+1}=\sum_{m=1}^{j}\theta_m$ y $R_{j+1}=\sum_{m=1}^{j}\delta_m$.

De \eqref{eq:delta_ki} se obtiene una relación recursiva explícita para $\theta_{j+1}$:
\begin{equation}
\boxed{
\theta_{j+1}
=
\frac{
2\cdot4^{N-j}\left[4k_j(\theta_j-\delta_j)+k_{j+1}(S_{j+1}-R_{j+1})+k_{j+1}\delta_{j+1}\right]
+
\frac{f}{2}L_{\text{quad}}\,\theta_j
}{
2\cdot4^{N-j}k_{j+1}
+
\frac{f}{2}L_{\text{quad}}
}.}
\label{eq:theta_recursive_ki}
\end{equation}

% =================================================
\section{Recurrencia adimensional y reducción a un problema efectivo}
\label{subsec:recursion_dimensionless}
% =================================================
A partir de la relación recursiva general \eqref{eq:theta_recursive_ki}, es útil trabajar con variables adimensionales. Definimos
\[
r_j := \frac{\theta_j}{\theta_1},
\qquad
s_j := \frac{1}{\theta_1}\sum_{m=1}^{j}\theta_m=\sum_{m=1}^{j}r_m,
\qquad
S_{j+1}= \sum_{m=1}^j \theta_m = s_j\,\theta_1.
\]

Para describir el estado inicial introducimos además
\[
d_j := \frac{\delta_j}{\theta_1},
\qquad
q_j := \frac{1}{\theta_1}\sum_{m=1}^{j}\delta_m
      = \sum_{m=1}^{j} d_m,
\qquad
R_{j+1} = q_j\,\theta_1 .
\]

Además, introducimos las cantidades
\begin{equation}
B_j := 2\cdot 4^{N-j},
\qquad
C := \frac{f}{2}\,L_{\text{quad}},
\label{eq:BjC_def}
\end{equation}
de modo que \eqref{eq:theta_recursive_ki} pueda escribirse de forma compacta como
\begin{equation}
\theta_{j+1}
=
\frac{
B_j\left(4k_j(\theta_j-\delta_j)+k_{j+1}(S_{j+1}-R_{j+1})+k_{j+1}\delta_{j+1}\right)
+C\,\theta_j
}{
B_j k_{j+1}+C
}.
\label{eq:theta_recursive_ki_BC}
\end{equation}

Dividiendo por $\theta_1$ y usando $\theta_j=r_j\theta_1$,
$S_{j+1}=s_j\theta_1$, $\delta_j=d_j\theta_1$ y $R_{j+1}=q_j\theta_1$,
obtenemos la recurrencia adimensional
\begin{equation}
\boxed{
r_{j+1}
=
\frac{
B_j\left[4k_j(r_j-d_j)+k_{j+1}(s_j-q_j)+k_{j+1}d_{j+1}\right]
+C\,r_j
}{
B_j k_{j+1}+C
},
}
\end{equation}
junto con las relaciones
\begin{equation}
s_{j+1}=s_j+r_{j+1},
\qquad
q_{j+1}=q_j+d_{j+1}.
\end{equation}

Las condiciones iniciales son
\[
r_1=1,
\qquad
s_1=1,
\qquad
q_1=d_1.
\]

\subsection{Forma tipo ``$\alpha_j$--$b_j$''}

La recurrencia puede escribirse en forma lineal introduciendo
\begin{equation}
b_j := \frac{B_j k_{j+1}}{B_j k_{j+1}+C},
\qquad
\alpha_j := \frac{4B_j k_j + C}{B_j k_{j+1}+C}.
\end{equation}

Con estas definiciones se obtiene
\begin{equation}
\boxed{
r_{j+1}=\alpha_j r_j + b_j s_j + \gamma_j,
\qquad
s_{j+1}=s_j+r_{j+1},
}
\label{eq:recurrencia}
\end{equation}
donde el término
\begin{equation}
\gamma_j :=
\frac{B_j}{B_jk_{j+1}+C}
\left[k_{j+1}(d_{j+1}-q_j)-4k_j d_j\right]
\end{equation}
representa el efecto del estado inicial.

\subsection{Reducción de la energía interna}

Usando $\theta_i=r_i\theta_1$ y $S_i=s_{i-1}\theta_1$, la parte elástica
del Hamiltoniano se reduce a
\begin{equation}
E_{\text{quad}}(\theta_1)
=
\sum_{i=1}^{N}4^{N-i+1}k_i
\left[(\theta_i-\delta_i)^2+(S_i-R_i)^2\right].
\end{equation}

Al sustituir las variables adimensionales obtenemos
\begin{equation}
E_{\text{quad}}(\theta_1)
=
\sum_{i=1}^{N}4^{N-i+1}k_i
\left[(r_i\theta_1-\delta_i)^2+(s_{i-1}\theta_1-R_i)^2\right].
\end{equation}

Expandiendo en potencias de $\theta_1$ se obtiene
\begin{equation}
E_{\text{quad}}(\theta_1)
=
A_k\,\theta_1^2
-
2B_\delta\,\theta_1
+
E_0,
\end{equation}
donde
\begin{align}
A_k &=
\sum_{i=1}^{N}4^{N-i+1}k_i\left(r_i^2+s_{i-1}^2\right),\\
B_\delta &=
\sum_{i=1}^{N}4^{N-i+1}k_i\left(r_i\delta_i+s_{i-1}R_i\right).
\end{align}

\subsection{Largo aproximado}

En el régimen de pequeños ángulos
\[
L_{\text{quad}}
\simeq
\exp\!\Big(\tfrac12\sum_{i=1}^N \theta_i
-
\tfrac14\sum_{i=1}^N \theta_i^2\Big),
\]
y al introducir $\theta_i=r_i\theta_1$ resulta
\begin{equation}
L_{\text{quad}}(\theta_1)
\simeq
\exp\!\Big(\tfrac{S}{2}\,\theta_1 - \tfrac{Q}{4}\,\theta_1^2\Big),
\qquad
S := \sum_{i=1}^{N} r_i,
\quad
Q := \sum_{i=1}^{N} r_i^{\,2}.
\end{equation}

\subsection{Hamiltoniano efectivo}

El Hamiltoniano efectivo queda
\begin{equation}
\boxed{
\mathcal{H}_{\text{quad}}(\theta_1)
=
A_k\,\theta_1^2
-
2B_\delta\,\theta_1
-
f\left[
\exp\!\Big(\tfrac{S}{2}\theta_1-\tfrac{Q}{4}\theta_1^2\Big)-\ell
\right].
}
\end{equation}

\subsection{Caso particular $\delta_i=0$}

El caso especial $\delta_i=0$ corresponde a un estado inicial no deformado.
En esta situación se tiene $R_i=0$ y el término lineal $B_\delta$ desaparece,
de modo que la energía elástica se reduce a
\begin{equation}
E_{\text{quad}}(\theta_1)=A_k\,\theta_1^2 .
\end{equation}

Además, la recurrencia adimensional \eqref{eq:recurrencia} se simplifica
a
\begin{equation}
r_{j+1}=\alpha_j r_j + b_j s_j,
\qquad
s_{j+1}=s_j+r_{j+1},
\end{equation}
la cual es lineal homogénea en $r_j$ y $s_j$. En consecuencia, los coeficientes
$\{r_i,s_i\}$ dependen únicamente de los parámetros del sistema
$\{k_i,B_i,C\}$ y no dependen de $\theta_1$ ni de ningún ángulo particular.
Esto permite calcular la secuencia $\{r_i\}$ de manera independiente y cerrar
el problema de forma analítica.

En particular, la condición de equilibrio para el último nivel,
Ec.~\eqref{eq:dH_thetaN_ki}, conduce a la relación
\begin{equation}
2\cdot4^{1}k_N\theta_N
-
\frac{f}{2}L\,
\frac{\cos(\tfrac{\theta_N}{2})-\sin(\tfrac{\theta_N}{2})}
{\cos(\tfrac{\theta_N}{2})+\sin(\tfrac{\theta_N}{2})}
=0 .
\end{equation}

En el régimen de pequeños ángulos,
\[
\frac{\cos(\tfrac{\theta_N}{2})-\sin(\tfrac{\theta_N}{2})}
{\cos(\tfrac{\theta_N}{2})+\sin(\tfrac{\theta_N}{2})}
\simeq
1-\theta_N+\frac{\theta_N^2}{2},
\]
lo que conduce a la ecuación cuadrática
\begin{equation}
\frac{fL}{2}\,\theta_N^2-(fL+16k_N)\theta_N+fL=0.
\end{equation}

Seleccionando la rama físicamente relevante (la solución pequeña), se obtiene
\begin{equation}
\theta_N=
\frac{(fL+16k_N)-\sqrt{(fL+16k_N)^2-2(fL)^2}}{fL}.
\end{equation}

Finalmente, usando la relación jerárquica $\theta_N=r_N\theta_1$, se obtiene
una expresión explícita para $\theta_1$:
\begin{equation}
\boxed{
\theta_1
=
\frac{1}{r_N}\,
\frac{(fL+16k_N)-\sqrt{(fL+16k_N)^2-2(fL)^2}}{fL}.
}
\end{equation}

% =================================================
\subsection{Análisis espectral}
\label{sec:espectral}
% =================================================

En esta sección reescribimos la ecuación de equilibrio en el espacio de Fourier con el objetivo de caracterizar la respuesta por modos. Para simplificar, fijaremos la constantes elásticas jerárquicas de la forma:
\begin{equation}
k_j = 4^{\,j-N},
\label{eq:kj_choice}
\end{equation}
la ecuación de equilibrio para $\Delta_j=0$ se simplifica, pues los prefactores jerárquicos se cancelan exactamente. En particular, se obtiene
\begin{equation}
D\,(\theta_{j+1}-\theta_j)
=
8\Big[(\delta_{j+1}-\delta_j) + (S_{j+1}-R_{j+1})\Big],
\qquad
D:=8+\frac{f}{2}L_{\rm quad},
\label{eq:theta_firstdiff_nonlocal}
\end{equation}
donde
\begin{equation}
S_{j+1}=\sum_{m=1}^{j}\theta_m,
\qquad
R_{j+1}=\sum_{m=1}^{j}\delta_m,
\label{eq:def_SR}
\end{equation}
son sumas acumuladas (no locales). La no localidad aparece porque la energía contiene términos del tipo $(S_i-R_i)^2$, de modo que la condición de equilibrio en un nivel $j$ ``siente'' la historia de ángulos en niveles anteriores.


\paragraph{Eliminación de la no localidad.} Una manera conveniente de eliminar explícitamente la suma acumulada consiste en considerar la diferencia discreta de \eqref{eq:theta_firstdiff_nonlocal} entre $\Delta_j$ y $\Delta_{j-1}$. Es decir, escribimos la misma ecuación para $j-1$,
\begin{equation}
D(\theta_{j}-\theta_{j-1})
=
8\Big[(\delta_{j}-\delta_{j-1})+(S_{j}-R_{j})\Big],
\label{eq:theta_firstdiff_nonlocal_jm1}
\end{equation}
y restamos \eqref{eq:theta_firstdiff_nonlocal_jm1} de \eqref{eq:theta_firstdiff_nonlocal}. Esto produce
\begin{equation}
D(\theta_{j+1}-2\theta_j+\theta_{j-1})
=
8\Big[(\delta_{j+1}-2\delta_j+\delta_{j-1}) + (S_{j+1}-R_{j+1})-(S_j-R_j)\Big].
\label{eq:second_diff_pre}
\end{equation}
Usando la identidad exacta
\begin{equation}
S_{j+1}=S_j+\theta_j,
\qquad
R_{j+1}=R_j+\delta_j
\;\;\Rightarrow\;\;
(S_{j+1}-R_{j+1})-(S_j-R_j)=\theta_j-\delta_j,
\label{eq:identity_S}
\end{equation}
la ecuación \eqref{eq:second_diff_pre} se vuelve completamente local:
\begin{equation}
D(\theta_{j+1}-2\theta_j+\theta_{j-1}) - 8\theta_j
=
8(\delta_{j+1}+\delta_{j-1}-3\delta_j).
\label{eq:local_helmholtz_discrete}
\end{equation}
La operación de diferencia discreta elimina la no localidad, pero introduce una constante de integración (modo global) que debe fijarse mediante condiciones de borde o imponiendo \eqref{eq:theta_firstdiff_nonlocal} en algún índice $j_\star$.

\paragraph{Convención de transformada discreta.}
Consideramos una cadena de longitud $M$ (índices $j=0,\dots,M-1$) con condiciones periódicas y usamos la transformada discreta de Fourier (DFT):
\begin{equation}
\hat\theta(k)=\sum_{j=0}^{M-1}\theta_j\,e^{-ikj},
\qquad
\theta_j=\frac{1}{M}\sum_{k}\hat\theta(k)e^{ikj},
\qquad
k=\frac{2\pi n}{M},\;n=0,\dots,M-1.
\label{eq:DFT_def}
\end{equation}
Bajo \eqref{eq:DFT_def}, desplazamientos discretos se diagonalizan como
\begin{equation}
\theta_{j\pm1} \;\xrightarrow{\mathcal F}\; e^{\pm ik}\hat\theta(k),
\qquad
\theta_{j+1}-2\theta_j+\theta_{j-1} \;\xrightarrow{\mathcal F}\; (e^{ik}-2+e^{-ik})\hat\theta(k)
= -4\sin^2\!\Big(\frac{k}{2}\Big)\hat\theta(k).
\label{eq:laplacian_Fourier}
\end{equation}
Análogamente,
\begin{equation}
\delta_{j+1}+\delta_{j-1}-3\delta_j
\;\xrightarrow{\mathcal F}\;
(e^{ik}+e^{-ik}-3)\hat\delta(k)
=
(2\cos k-3)\hat\delta(k).
\label{eq:source_Fourier}
\end{equation}

\paragraph{Relación espectral $\hat\theta(k)$.}
Aplicando \eqref{eq:laplacian_Fourier}--\eqref{eq:source_Fourier} a \eqref{eq:local_helmholtz_discrete} obtenemos, para cada modo $k$,
\begin{equation}
\Big[-4D\sin^2\!\Big(\frac{k}{2}\Big)-8\Big]\hat\theta(k)
=
8(2\cos k-3)\hat\delta(k),
\label{eq:theta_k_equation}
\end{equation}
y por tanto
\begin{equation}
\boxed{
\hat\theta(k)=
-\frac{8(2\cos k-3)}{4D\sin^2(k/2)+8}\,\hat\delta(k),
\qquad
D=8+\frac{f}{2}L_{\rm quad}.
}
\label{eq:theta_k_solution}
\end{equation}
La ecuación \eqref{eq:theta_k_solution} muestra que cada modo espectral responde mediante un ``filtro'' racional en $\sin^2(k/2)$, con una amplitud controlada por la fuerza $f$ a través de $D$. El modo $k=0$ requiere, adicionalmente, consistencia con la condición global eliminada al pasar de \eqref{eq:theta_firstdiff_nonlocal} a \eqref{eq:local_helmholtz_discrete}.

% -----------------------------
\bibliographystyle{unsrtnat}
\bibliography{refs}
% -----------------------------

\end{document}
